\setcounter{table}{0}
\renewcommand{\thetable}{D.\arabic{table}}
\renewcommand{\theHtable}{D.\arabic{table}}

Este anexo describe el procedimiento necesario para compilar, cargar y poner en
funcionamiento el firmware de bajo nivel sobre el hexápodo, así como para instalar,
compilar y lanzar el subsistema ROS~2 que se ejecuta en la Raspberry~Pi. El objetivo
no es repetir la arquitectura desarrollada en el Capítulo~4, sino proporcionar una
secuencia reproducible de preparación, arranque, comprobación y parada para ambos
subsistemas.

\section*{D.1. Requisitos previos y seguridad}

Para cargar el firmware se necesita un equipo con ESP-IDF~5.5.1, un cable USB
de programación y acceso al directorio del proyecto \texttt{firmware\_esp32}. En el
robot, los 18 AX-12A deben estar configurados a 1~Mbit/s y con los
identificadores indicados en la Tabla~\ref{tab:manual_ids_servos}. La PCB debe
proporcionar la adaptación \textit{half-duplex} y la alimentación descritas en
las Subsecciones~\ref{subsec:hw_half_duplex} y~\ref{subsec:driver_ax12}.

\noindent\fcolorbox{black}{gray!10}{%
    \parbox{0.94\textwidth}{%
        \textbf{Precaución.} Durante la primera puesta en marcha el hexápodo
        debe permanecer elevado, con todas las patas libres y sin personas u
        objetos dentro de su recorrido. Al arrancar, el firmware activa el par
        de los actuadores y envía la postura neutral, aunque el modo de marcha
        inicial sea \texttt{STOP}. Todas las conexiones deben realizarse con la
        alimentación de los actuadores desconectada.
    }%
}

\begin{table}[H]
    \centering
    \small
    \begin{tabularx}{\textwidth}{@{}Xccc@{}}
        \toprule
        \textbf{Pata} & \textbf{Coxa} & \textbf{Fémur} & \textbf{Tibia} \\
        \midrule
        Izquierda delantera & 1  & 3  & 5  \\
        Izquierda media     & 13 & 15 & 17 \\
        Izquierda trasera   & 7  & 9  & 11 \\
        Derecha trasera     & 8  & 10 & 12 \\
        Derecha media       & 14 & 16 & 18 \\
        Derecha delantera   & 2  & 4  & 6  \\
        \bottomrule
    \end{tabularx}
    \caption{Identificadores de los AX-12A utilizados por el firmware.}
    \label{tab:manual_ids_servos}
\end{table}

\section*{D.2. Conexiones utilizadas por el firmware}

El ESP32 emplea tres UART diferentes. La UART0 queda reservada para la carga y
los mensajes de depuración; la UART1 recibe los comandos de control; y la UART2
se conecta al circuito de adaptación del bus DYNAMIXEL. La configuración
utilizada se resume en la Tabla~\ref{tab:manual_uart}.

\begin{table}[H]
    \centering
    \small
    \begin{tabularx}{\textwidth}{@{}lllX@{}}
        \toprule
        \textbf{Interfaz} & \textbf{Señales} & \textbf{Velocidad} & \textbf{Uso} \\
        \midrule
        UART0 & USB/programación & 115200~bit/s & Carga del firmware y registro de arranque. \\
        UART1 & TX: GPIO4; RX: GPIO5 & 115200~bit/s & Comandos y telemetría. \\
        UART2 & TX: GPIO16; RX: GPIO17; DIR: GPIO2 & 1~Mbit/s & Bus \textit{half-duplex} de los AX-12A. \\
        \bottomrule
    \end{tabularx}
    \caption{Interfaces serie empleadas en el ESP32.}
    \label{tab:manual_uart}
\end{table}

Para una prueba directa de la UART1 puede utilizarse un adaptador USB--UART de
3,3~V. Su RX debe conectarse al GPIO4 del ESP32, su TX al GPIO5 y ambas placas
deben compartir GND. No debe aplicarse una señal serie de 5~V al ESP32. Cada
comando se transmite como una línea de texto terminada en salto de línea. La
consola USB de programación no sustituye esta conexión, ya que corresponde a la
UART0.

\section*{D.3. Compilación y carga}

Con el entorno de ESP-IDF inicializado, el firmware se compila y carga mediante:

\begin{lstlisting}[basicstyle=\ttfamily\small,frame=single,breaklines=true]
cd /ruta/al/proyecto/firmware_esp32
idf.py build
idf.py -p /dev/ttyUSB0 flash monitor
\end{lstlisting}

El nombre del puerto puede ser distinto, por ejemplo \texttt{/dev/ttyACM0}. La
orden \texttt{monitor} muestra la salida de depuración de la UART0 y se abandona
con \texttt{Ctrl+]}.

La carga puede realizarse sin alimentar los actuadores. Sin embargo, antes de
la puesta en marcha real debe estar disponible el bus completo y debe
reiniciarse el ESP32, puesto que el escaneo y la activación de par se ejecutan
una sola vez durante el arranque.

\section*{D.4. Primera puesta en marcha}

Con el robot elevado y las conexiones revisadas, el procedimiento recomendado
es el siguiente:

\begin{enumerate}[leftmargin=*,label=\arabic*.]
    \item Mantener inicialmente desconectada la alimentación de potencia y
          comprobar la polaridad, la masa común y los seis ramales del bus.
    \item Conectar la interfaz de depuración y, si se desea operar sin ROS~2,
          un terminal serie independiente a la UART1 a 115200~bit/s, 8N1.
    \item Alimentar el sistema o reiniciar el ESP32 con los AX-12A ya
          disponibles en el bus.
    \item Verificar en la consola que el escaneo detecta los servos con IDs
          1--18 y finaliza con el mensaje \texttt{18/18 respuestas}.
    \item Confirmar que las patas adoptan la postura neutral sin colisiones,
          vibraciones ni movimientos en sentido contrario al esperado.
\end{enumerate}

El estado inicial del firmware es \texttt{STOP}, con periodo de marcha de 2~s, altura de
paso de 40~mm, postura local nula y tiempo de suavizado de 1~s. Cuando se
utiliza ROS~2, el puente modifica esta configuración al conectarse
(Sección~D.12). El modo
\texttt{STOP} evita el avance, pero mantiene el par de los actuadores.

\section*{D.5. Prueba funcional mediante UART1}

Antes de apoyar el robot en el suelo debe realizarse una prueba a baja
velocidad. La siguiente secuencia restablece una configuración conocida, activa
el patrón trípode y solicita un avance longitudinal reducido:

\begin{lstlisting}[basicstyle=\ttfamily\small,frame=single]
HELP
MOD,0
CMD,0,0,0
POS,0,0,0,0,0,0
SMO,2
GAT,2,20
MOD,1
CMD,20,0,0
\end{lstlisting}

Tras comprobar varios ciclos, la parada se realiza enviando primero
\texttt{CMD,0,0,0}. Cuando haya cesado el movimiento, se envía \texttt{MOD,0}.
Las consignas permanecen almacenadas hasta recibir una nueva orden o reiniciar
el sistema; por ello no debe desconectarse el enlace de control mientras exista
una velocidad distinta de cero.

\begin{table}[H]
    \centering
    \footnotesize
    \begin{tabularx}{\textwidth}{@{}>{\ttfamily}lX@{}}
        \toprule
        \normalfont\textbf{Comando} & \textbf{Función} \\
        \midrule
        HELP o ? & Muestra la ayuda integrada. \\
        MOD,m & Selecciona \texttt{0=STOP}, \texttt{1=TRIPOD}, \texttt{2=RIPPLE} o \texttt{3=AMBLE}. \\
        CMD,vx,vy,vyaw & Fija las velocidades $v_x$ y $v_y$ en mm/s y $v_{yaw}$ en rad/s. \\
        POS,r,p,y,x,ty,z & Fija la postura: rotaciones en radianes y desplazamientos en milímetros. \\
        GAT,T,H & Fija el periodo $T$ en segundos y la altura $H$ del paso en milímetros. \\
        SMO,t & Fija el tiempo de suavizado de la postura en segundos. \\
        RST & Reinicia el origen de la odometría. \\
        REG,id,a,n & Lee uno o dos bytes de la tabla de control de un AX-12A. \\
        WREG,id,a,n,v & Escribe uno o dos bytes en la tabla de control de un AX-12A. \\
        \bottomrule
    \end{tabularx}
    \caption{Comandos aceptados por la UART1.}
    \label{tab:manual_comandos_uart}
\end{table}

El controlador limita individualmente \textit{roll} y \textit{pitch} a
$\pm0{,}35$~rad, el \textit{yaw} local a $\pm0{,}52$~rad, los desplazamientos
$x$ e $y$ a $\pm40$~mm y $z$ a $\pm50$~mm. El periodo de marcha no puede ser
inferior a 0,6~s y la altura de paso se mantiene entre 10 y 60~mm. Además, el
generador normaliza conjuntamente el desplazamiento efectivo de cada paso, es
decir, el que recorre el cuerpo durante una fase de apoyo, con límites de
120~mm en X, 50~mm en Y y 0,35~rad de giro. Estos límites no evitan
que una combinación extrema de postura y marcha pueda producir una cinemática
inversa inválida, por lo que la validación inicial debe hacerse con consignas
reducidas.

\noindent\textbf{Advertencia sobre \texttt{WREG}.} Este comando permite
modificar cualquier registro admitido por el firmware, incluidos parámetros de
EEPROM como el identificador o la velocidad del bus. Debe reservarse para
diagnóstico y utilizarse solo conociendo la tabla de control del AX-12A.

\section*{D.6. Telemetría y comprobación}

La UART1 transmite automáticamente una línea \texttt{ODO} cada segundo y
publica los resultados del diagnóstico de los actuadores. Las respuestas se
interpretan de acuerdo con la Tabla~\ref{tab:manual_respuestas_uart}.

\begin{table}[H]
    \centering
    \footnotesize
    \begin{tabularx}{\textwidth}{@{}>{\ttfamily}lX@{}}
        \toprule
        \normalfont\textbf{Respuesta} & \textbf{Contenido} \\
        \midrule
        ODO,x,y,yaw,roll,pitch,body\_yaw,x\_body,y\_body,z\_body & Odometría plana y postura local confirmada. \\
        AXH,id,valid,V,T,load,error & Validez, tensión, temperatura, carga y byte de error de un actuador. La tensión se expresa en décimas de voltio. \\
        AXR,id,a,n,valid,value,error & Resultado de una lectura solicitada con \texttt{REG}. \\
        AXW,id,a,n,valid,value,error & Resultado de una escritura solicitada con \texttt{WREG}. \\
        \bottomrule
    \end{tabularx}
    \caption{Mensajes transmitidos por el firmware.}
    \label{tab:manual_respuestas_uart}
\end{table}

Una muestra \texttt{AXH} no válida o con un byte de error distinto de cero
debe investigarse antes de apoyar el robot. La odometría publicada es una
estimación obtenida a partir de la consigna de marcha y no sustituye una medida
externa del desplazamiento real.

\section*{D.7. Parada y resolución de problemas}

La parada ordinaria consiste en enviar \texttt{CMD,0,0,0}, esperar a que las
patas terminen la transición y seleccionar \texttt{MOD,0}. Esta orden no
desactiva el par; para manipular las articulaciones manualmente debe
desconectarse la alimentación de los actuadores. Si se pierde la comunicación
con una consigna de velocidad activa, debe cortarse dicha alimentación.

Los fallos más habituales pueden localizarse con las siguientes comprobaciones:

\begin{itemize}[leftmargin=*,nosep]
    \small
    \item \textbf{No se detecta ningún servo:} comprobar la alimentación del
          bus, la masa común, la adaptación \textit{half-duplex}, la velocidad
          de 1~Mbit/s y las señales GPIO16, GPIO17 y GPIO2.
    \item \textbf{El escaneo es incompleto:} revisar el ID, el cableado y la
          alimentación del ramal correspondiente. No debe haber IDs repetidos.
    \item \textbf{No se reciben comandos:} comprobar que se utiliza la UART1,
          que TX y RX están cruzados, que el formato es 115200~bit/s 8N1 y que
          cada orden termina en salto de línea.
    \item \textbf{Se registra una IK inválida:} enviar una velocidad nula,
          recuperar la postura neutral y reducir la altura, la velocidad o el
          desplazamiento solicitado antes de reanudar la marcha.
    \item \textbf{Una articulación gira en sentido incorrecto:} detener el
          sistema y revisar la correspondencia de IDs, los signos y los ajustes
          de montaje definidos en \texttt{leg\_config.cpp}.
\end{itemize}

\section*{D.8. Requisitos e instalación de ROS~2 en la Raspberry~Pi}

El subsistema de alto nivel, descrito en la Sección~\ref{sec:integracion_ros2}, se
ejecuta sobre una Raspberry~Pi~4 \cite{raspberrypi4_datasheet} con Ubuntu Server
24.04~LTS de 64~bits, la distribución de base admitida por ROS~2 Jazzy. Sobre esa
base se instalan el paquete \texttt{ros-jazzy-ros-base}, la herramienta de
compilación \texttt{colcon} (\texttt{python3-colcon-common-extensions}) y el
gestor de dependencias \texttt{rosdep} (\texttt{python3-rosdep}), siguiendo el
procedimiento oficial de instalación de ROS~2 \cite{ros2_docs}. Instalado el
entorno, cada sesión debe cargarlo:

\begin{lstlisting}[basicstyle=\ttfamily\small,frame=single]
source /opt/ros/jazzy/setup.bash
\end{lstlisting}

por lo que conviene añadir esta línea a \texttt{\textasciitilde/.bashrc}. La
Tabla~\ref{tab:manual_deps_ros2} resume las dependencias declaradas por cada
paquete del proyecto, que \texttt{rosdep} resuelve automáticamente en la
Sección~D.11.

\begin{table}[H]
    \centering
    \small
    \begin{tabularx}{\textwidth}{@{}l l X@{}}
        \toprule
        \textbf{Paquete} & \textbf{Compilación} & \textbf{Dependencias} \\
        \midrule
        \texttt{hexapod\_bridge} & \texttt{ament\_python} & \texttt{rclpy}, \texttt{geometry\_msgs}, \texttt{nav\_msgs}, \texttt{std\_msgs}, \texttt{tf\_transformations}, \texttt{python3-serial} \\
        \texttt{hexapod\_web} & \texttt{ament\_python} & \texttt{rclpy}, \texttt{geometry\_msgs}, \texttt{nav\_msgs}, \texttt{std\_msgs}, \texttt{tf\_transformations}, \texttt{python3-flask} \\
        \texttt{obstacle\_avoider} & \texttt{ament\_python} & \texttt{rclpy}, \texttt{sensor\_msgs}, \texttt{geometry\_msgs}, \texttt{python3-numpy} \\
        \texttt{sllidar\_ros2} & \texttt{ament\_cmake} & \texttt{rclcpp}, \texttt{std\_srvs}, \texttt{sensor\_msgs} \\
        \texttt{hexapod\_bringup} & \texttt{ament\_cmake} & \texttt{launch}, \texttt{launch\_ros}, \texttt{ament\_index\_python}, \texttt{laser\_filters} y los cuatro paquetes anteriores \\
        \bottomrule
    \end{tabularx}
    \caption{Dependencias declaradas por los paquetes ROS~2 del hexápodo.}
    \label{tab:manual_deps_ros2}
\end{table}

\section*{D.9. Enlace serie con el ESP32}

A diferencia de la prueba directa descrita en la Sección~D.2, en funcionamiento
normal el puente \texttt{/esp32\_bridge} se comunica con el ESP32 a través del
UART por hardware de la propia Raspberry~Pi (\texttt{/dev/ttyAMA0}), sin
adaptador USB--UART intermedio: el GPIO14 (TX) se conecta al GPIO5 (RX) del
ESP32, el GPIO15 (RX) al GPIO4 (TX) del ESP32, y ambas placas comparten masa.
Al ser ambos extremos de 3,3~V, no se necesita la adaptación de niveles que sí
exige el bus de los AX-12A.

En la Raspberry~Pi~4, la interfaz PL011 que corresponde a \texttt{/dev/ttyAMA0}
está asignada por defecto al controlador Bluetooth, por lo que debe liberarse y
habilitarse manualmente. En \texttt{/boot/firmware/config.txt} se añade:

\begin{lstlisting}[basicstyle=\ttfamily\small,frame=single]
enable_uart=1
dtoverlay=disable-bt
\end{lstlisting}

y se deshabilitan los servicios que, de lo contrario, reutilizarían el puerto
como consola serie:

\begin{lstlisting}[basicstyle=\ttfamily\small,frame=single,breaklines=true]
sudo systemctl disable --now hciuart.service
sudo systemctl disable --now serial-getty@ttyAMA0.service
sudo usermod -aG dialout $USER
\end{lstlisting}

Tras reiniciar, \texttt{/dev/ttyAMA0} debe existir y ser accesible sin
privilegios de superusuario; este es el puerto y la velocidad (115200~bit/s)
configurados por defecto en \texttt{hexapod\_bridge/config/esp32\_bridge.yaml}.

\noindent\fcolorbox{black}{gray!10}{%
    \parbox{0.94\textwidth}{%
        \textbf{Precaución.} Un error en \texttt{config.txt} puede impedir el
        arranque gráfico o dejar la Raspberry~Pi sin consola serie de
        emergencia. Conviene conservar una copia del archivo original y
        disponer de acceso por SSH antes de reiniciar.
    }%
}

\section*{D.10. Preparación del LiDAR}

El RPLIDAR~A1M8 se conecta por USB y se presenta como un adaptador serie
CP210x. Para evitar que el nombre del dispositivo cambie según el orden de
conexión, \texttt{sllidar\_ros2} \cite{sllidar_ros2} incluye la siguiente
regla \texttt{udev}, que debe copiarse a \texttt{/etc/udev/rules.d/} y
recargarse:

\begin{lstlisting}[basicstyle=\ttfamily\small,frame=single,breaklines=true]
KERNEL=="ttyUSB*", ATTRS{idVendor}=="10c4", ATTRS{idProduct}=="ea60", MODE:="0777", SYMLINK+="rplidar"
\end{lstlisting}

\begin{lstlisting}[basicstyle=\ttfamily\small,frame=single,breaklines=true]
sudo cp src/sllidar_ros2/scripts/rplidar.rules /etc/udev/rules.d/
sudo udevadm control --reload-rules
sudo udevadm trigger
\end{lstlisting}

A partir de ese momento el sensor queda disponible de forma estable en
\texttt{/dev/rplidar}, con independencia del puerto \texttt{ttyUSB} que le
asigne el sistema, y es ese enlace el que conviene indicar al lanzamiento
descrito en la Sección~D.12.

\section*{D.11. Obtención y compilación del espacio de trabajo}

Con ROS~2 instalado y \texttt{colcon} disponible, el código del proyecto
\cite{repositorio_hexapodo} se sitúa dentro de \texttt{src} de un espacio de
trabajo y se compila:

\begin{lstlisting}[basicstyle=\ttfamily\small,frame=single,breaklines=true]
mkdir -p ~/ros2_ws/src
git clone https://github.com/josuebcc/hexapod-ros2 ~/ros2_ws/src
cd ~/ros2_ws
rosdep install --from-paths src --ignore-src -r -y
colcon build --symlink-install
source install/setup.bash
\end{lstlisting}

El espacio de trabajo contiene cinco paquetes: los tres desarrollados para
este proyecto (\texttt{hexapod\_bridge}, \texttt{hexapod\_web} y
\texttt{obstacle\_avoider}), el controlador del fabricante \texttt{sllidar\_ros2}
\cite{sllidar_ros2} y el paquete de lanzamiento \texttt{hexapod\_bringup}, que
no contiene código propio más allá del archivo de lanzamiento descrito a
continuación. Como con el firmware, \texttt{source install/setup.bash} debe
añadirse a \texttt{\textasciitilde/.bashrc} para no repetirlo en cada sesión.

\section*{D.12. Lanzamiento del sistema}

El lanzamiento conjunto se realiza con un único comando:

\begin{lstlisting}[basicstyle=\ttfamily\small,frame=single]
ros2 launch hexapod_bringup hexapod.launch.py
\end{lstlisting}

que arranca \texttt{/esp32\_bridge}, \texttt{/web\_panel}, \texttt{/reactive\_node}
(el nodo de \texttt{obstacle\_avoider}), el filtro
\texttt{/scan\_to\_scan\_filter\_chain} de \texttt{laser\_filters} y el controlador
\texttt{sllidar\_ros2}, en las condiciones resumidas en la
Tabla~\ref{tab:manual_launch_args}. El puente no
incorpora valores por defecto: si se ejecuta de forma aislada con
\texttt{ros2 run}, debe recibir siempre el archivo \texttt{esp32\_bridge.yaml}.

Unos dos segundos después de arrancar, \texttt{/esp32\_bridge} envía una única
vez la configuración inicial definida en \texttt{esp32\_bridge.yaml}. Con los
valores por defecto, esta selecciona la marcha trípode (\texttt{MOD,1}), un
periodo de 2~s con altura de paso de 30~mm (\texttt{GAT,2,30}) y un suavizado
de 0,5~s (\texttt{SMO,0.5}). Por tanto, a diferencia del arranque descrito en la
Sección~D.4, el robot no permanece en \texttt{STOP}, aunque no avanza mientras
no reciba una velocidad en \path{/hexapod/cmd_vel}.

\begin{table}[H]
    \centering
    \footnotesize
    \begin{tabularx}{\textwidth}{@{}l l X@{}}
        \toprule
        \textbf{Argumento} & \textbf{Valor por defecto} & \textbf{Función} \\
        \midrule
        \texttt{lidar\_port} & \texttt{/dev/ttyUSB0} & Puerto serie del LiDAR (usar \texttt{/dev/rplidar} tras la Sección~D.10). \\
        \texttt{lidar\_baud} & \texttt{115200} & Velocidad de transmisión del LiDAR. \\
        \texttt{start\_web} & \texttt{true} & Arranca o no el panel web. \\
        \texttt{start\_lidar} & \texttt{true} & Arranca o no el controlador del LiDAR. \\
        \texttt{bridge\_params\_file} & \texttt{esp32\_bridge.yaml} & Parámetros del puente ESP32. \\
        \texttt{web\_params\_file} & \texttt{web\_panel.yaml} & Parámetros del panel web. \\
        \texttt{scan\_filters\_file} & \texttt{scan\_filters.yaml} & Filtros aplicados al barrido del LiDAR. \\
        \bottomrule
    \end{tabularx}
    \caption{Argumentos de \texttt{hexapod.launch.py}.}
    \label{tab:manual_launch_args}
\end{table}

Por ejemplo, para lanzar el sistema con el LiDAR fijado por la regla
\texttt{udev} y sin el panel web:

\begin{lstlisting}[basicstyle=\ttfamily\small,frame=single,breaklines=true]
ros2 launch hexapod_bringup hexapod.launch.py lidar_port:=/dev/rplidar start_web:=false
\end{lstlisting}

El evitador también puede ejecutarse de forma aislada, con
\texttt{/scan\_filtered} como entrada y \texttt{/obstacle\_avoider/cmd\_vel} como
salida, siempre que el filtro del LiDAR esté en ejecución:

\begin{lstlisting}[basicstyle=\ttfamily\small,frame=single]
ros2 run obstacle_avoider reactive_node
\end{lstlisting}

\section*{D.13. Uso del panel web}

Con \texttt{/web\_panel} en ejecución, la interfaz queda disponible en el
puerto~8080 de la Raspberry~Pi, por ejemplo en
\url{http://192.168.1.145:8080/} dentro de la red del robot. El panel es el
único nodo que publica en \path{/hexapod/cmd_vel} y arranca siempre en modo
manual: en ese modo reenvía sus propios controles, mientras que en modo
autónomo reenvía las propuestas recibidas en \path{/obstacle_avoider/cmd_vel}. La
activación del modo autónomo, mediante las rutas \texttt{GET}/\texttt{POST}
\texttt{/control\_mode}, exige una propuesta del evitador recibida hace menos
de 0,5~s; en caso contrario se rechaza. Durante la autonomía el servidor
descarta además cualquier velocidad manual con el código HTTP~409, y una
parada (\texttt{MOD,0}) o el cierre del nodo devuelven el sistema a modo
manual. La Tabla~\ref{tab:manual_limites_postura} recoge los límites que el
panel aplica a la postura del cuerpo, configurables en
\texttt{web\_panel.yaml}.

\begin{table}[H]
    \centering
    \small
    \begin{tabularx}{\textwidth}{@{}l c c c@{}}
        \toprule
        \textbf{Parámetro} & \textbf{Mínimo} & \textbf{Máximo} & \textbf{Paso} \\
        \midrule
        Roll & $-20^\circ$ & $20^\circ$ & $1^\circ$ \\
        Pitch & $-20^\circ$ & $20^\circ$ & $1^\circ$ \\
        Yaw local & $-25^\circ$ & $25^\circ$ & $1^\circ$ \\
        $x$ & $-40$~mm & $40$~mm & $1$~mm \\
        $y$ & $-40$~mm & $40$~mm & $1$~mm \\
        $z$ & $-50$~mm & $50$~mm & $1$~mm \\
        \bottomrule
    \end{tabularx}
    \caption{Límites de postura configurados en \texttt{web\_panel.yaml}.}
    \label{tab:manual_limites_postura}
\end{table}

La vista de servos de la misma interfaz permite leer y escribir registros
individuales de los AX-12A mediante los comandos \texttt{REG}/\texttt{WREG}
descritos en la Sección~D.5; la escritura no vuelve a leer el registro de
forma automática, por lo que la comprobación debe hacerse manualmente con el
botón \textit{Leer}.

\section*{D.14. Resolución de problemas (Raspberry~Pi y ROS~2)}

\begin{itemize}[leftmargin=*,nosep]
    \small
    \item \textbf{\texttt{ros2 launch} no encuentra los paquetes:} falta
          cargar \texttt{/opt/ros/jazzy/setup.bash} o
          \texttt{install/setup.bash}, o \texttt{colcon build} terminó con
          errores que no se revisaron.
    \item \textbf{Permiso denegado en \texttt{/dev/ttyAMA0},
          \texttt{/dev/ttyUSB0} o \texttt{/dev/rplidar}:} el usuario no
          pertenece al grupo \texttt{dialout}; tras añadirlo es necesario
          cerrar sesión y volver a entrar.
    \item \textbf{El puente no recibe telemetría del ESP32:} el UART de la
          Raspberry~Pi no está habilitado, el servicio
          \texttt{serial-getty@ttyAMA0} sigue activo, o el puerto y la
          velocidad de \texttt{esp32\_bridge.yaml} no coinciden con los de la
          Sección~D.2.
    \item \textbf{El LiDAR no aparece o cambia de nombre:} falta instalar la
          regla \texttt{udev} de la Sección~D.10; debe usarse
          \texttt{/dev/rplidar} en lugar de \texttt{/dev/ttyUSBn}.
    \item \textbf{El panel web no responde:} comprobar que ningún otro
          proceso ocupa el puerto~8080 y que la Raspberry~Pi es accesible en
          la red utilizada.
    \item \textbf{No es posible activar el modo autónomo:}
          \texttt{reactive\_node} no está en ejecución o no recibe
          \texttt{/scan\_filtered} (\texttt{start\_lidar} desactivado, LiDAR
          desconectado o filtro detenido), o su última propuesta supera los
          0,5~s exigidos.
    \item \textbf{El robot no se detiene al perder la conexión:} verificar
          que \texttt{/esp32\_bridge} sigue en ejecución y que
          \texttt{cmd\_vel\_timeout\_s} en \texttt{esp32\_bridge.yaml}
          conserva un valor bajo (1~s por defecto).
\end{itemize}
